\documentclass[11pt]{amsart}
\usepackage{amsmath,amssymb,amsthm,mathtools}
\usepackage[margin=1.1in]{geometry}
\usepackage{booktabs}
\usepackage[no-math]{fontspec}
\newfontfamily\leanfont{Menlo}[Scale=0.8]
\usepackage{fancyvrb}
\usepackage{xcolor}
\usepackage{hyperref}
\hypersetup{colorlinks=true,linkcolor=blue!50!black,citecolor=blue!50!black,urlcolor=blue!50!black}

\newtheorem{theorem}{Theorem}[section]
\newtheorem{lemma}[theorem]{Lemma}
\newtheorem{proposition}[theorem]{Proposition}
\newtheorem{corollary}[theorem]{Corollary}
\theoremstyle{definition}
\newtheorem{definition}[theorem]{Definition}
\newtheorem{remark}[theorem]{Remark}

\newcommand{\R}{\mathbb{R}}
\newcommand{\C}{\mathbb{C}}
\newcommand{\Q}{\mathbb{Q}}
\newcommand{\Z}{\mathbb{Z}}
\newcommand{\cH}{\mathcal{H}}
\newcommand{\E}{\mathbb{E}}
\newcommand{\PP}{\mathbb{P}}
\DeclareMathOperator{\Haar}{Haar}

\title[Arbitrary and measurable colorings of Euclidean space]{Arbitrary colorings of Euclidean space are no easier than measurable ones}
\author{Henry Bass}
\date{October 2026}



\begin{document}
\begin{abstract}
A recent preprint released by OpenAI proves that every proper $k$-coloring of the plane, with arbitrary color classes, can be converted into a measurable $k$-coloring whose monochromatic unit pairs form a null set. It uses this to show that the plane is not $5$-colorable. Its averaging step relies on the rotation group of the plane being abelian, and experts have noted that the corresponding step fails in dimension $3$ and higher because the rotation group is not amenable.

We extend the transfer to every dimension $n\ge2$. We average only over independent rotations of orthogonal coordinate $2$-planes. We apply the planar rigidity theorem block by block through a polarization identity, and in odd dimensions we remove the remaining axis with an invariant-mean line average. Consequently
\[
\chi(\R^n)\ \ge\ \bigl\lceil 1/m_1(\R^n)\bigr\rceil ,
\]
where $m_1(\R^n)$ is the maximal density of a measurable set avoiding distance $1$. With published bounds on $m_1$ this gives $\chi(\R^3)\ge7$ (previously $6$), $\chi(\R^4)\ge 11$ (previously $9$), $\chi(\R^8)\ge53$ (previously $19$), and $\chi(\R^n)\ge(1.268+o(1))^n$, improving Raigorodskii's bound $(1.239+o(1))^n$ from 2000.

The three-dimensional transfer and the deduction of $\chi(\R^3)\ge7$ from the published density bound are formally verified in Lean~4. Applying the same argument to random independent sets, we show that finite unit-distance graphs attain the measurable density in every dimension, $\inf_G \alpha(G)/|G| = m_1(\R^n)$. In the plane this proves the first part of a recent conjecture of D\'ucz and Varga: Erd\H{o}s's function satisfies $f(N)=(m_1(\R^2)+o(1))N$.
\end{abstract}
\maketitle

\section{Introduction}

A proper $k$-coloring of $\R^n$ is a map $c\colon\R^n\to\{1,\dots,k\}$ with $c(x)\neq c(y)$ whenever $|x-y|=1$. The chromatic number $\chi(\R^n)$ is the least such $k$, with no regularity imposed on the color classes. For the plane, Nelson and Isbell showed $4\le\chi(\R^2)\le7$ around 1950. De Grey raised the lower bound to $5$ in 2018~\cite{deGrey}. A preprint released by OpenAI in September 2026~\cite{OAI} proves $\chi(\R^2)\ge6$, with a Lean formalization that we have rebuilt and checked. In three dimensions the best previous bounds were $6\le\chi(\R^3)\le15$~\cite{Nechushtan,Coulson,RT}. Asymptotically, $(1.239+o(1))^n\le\chi(\R^n)\le(3+o(1))^n$~\cite{Raigorodskii,LR}.

The measurable chromatic number $\chi_m(\R^n)$ requires measurable color classes. It is bounded below by $1/m_1(\R^n)$. Here
\[
m_1(\R^n)=\sup\{\overline{\delta}(A): A\subseteq\R^n\text{ measurable, no two points of }A\text{ at distance }1\},
\]
and $\overline\delta$ denotes upper density with respect to centered cubes. Fourier-analytic and semidefinite methods~\cite{OV,BPT,DOV} give much stronger lower bounds for $\chi_m$ than combinatorial methods give for $\chi$. For example, $\chi_m(\R^3)\ge7$, and $\chi_m(\R^n)\ge(1.268+o(1))^n$.

The OpenAI argument has two parts. The first is a transfer theorem~\cite[Theorem 1.3]{OAI}: the plane has a proper $k$-coloring if and only if it has a \emph{weak measurable} $k$-coloring. The second shows that weak measurable $5$-colorings do not exist. The transfer averages a coloring of the algebraic points over algebraic translations and rotations. This needs the group $\R^2_{\mathrm{alg}}\rtimes SO(2)_{\mathrm{alg}}$ to be amenable, which holds because it is solvable. Matolcsi, Ruzsa, Varga and Zsámboki~\cite{MRVZ} use the same kind of averaging for fractional colorings. They point out that it ``cannot automatically be extended to higher dimensions,'' since the Følner property fails for Euclidean motions when $d\ge3$.

\begin{definition}\label{def:weak}
A \emph{weak measurable $k$-coloring} of $\R^n$ is a Lebesgue measurable $c\colon\R^n\to\{1,\dots,k\}$ with classes $A_i=c^{-1}(i)$ such that for every $R>0$
\[
\sum_{i=1}^k\int_{B(0,R)}\int_{S^{n-1}}\mathbf 1_{A_i}(x)\,\mathbf 1_{A_i}(x+u)\,d\sigma(u)\,dx=0 ,
\]
where $\sigma$ is normalized surface measure on the unit sphere.
\end{definition}

\begin{theorem}[Transfer in every dimension]\label{thm:A}
Let $n\ge2$ and $k\ge1$. If $\R^n$ has a proper $k$-coloring, then it has a weak measurable $k$-coloring.
\end{theorem}

The converse holds by the density-point argument of~\cite[Lemma 4.3 and proof of Theorem 1.3]{OAI}, which is dimension-free. Thus $\chi(\R^n)$ equals the weak measurable chromatic number for every $n$.

\begin{corollary}\label{cor:B}
If $\R^n$ has a proper $k$-coloring, then there are pairwise disjoint Borel sets $A_1^*,\dots,A_k^*\subseteq\R^n$, whose union has null complement, such that no $A_i^*$ contains two points at distance $1$. Consequently
\[
\chi(\R^n)\ \ge\ \lceil 1/m_1(\R^n)\rceil .
\]
\end{corollary}

\begin{corollary}\label{cor:C}
$\chi(\R^3)\ge7$, $\chi(\R^4)\ge11$, $\chi(\R^5)\ge17$, $\chi(\R^6)\ge23$, $\chi(\R^7)\ge39$, and $\chi(\R^8)\ge53$. Moreover
\[
\chi(\R^n)\ \ge\ (1.268+o(1))^n\qquad(n\to\infty).
\]
\end{corollary}

The bounds for $3\le n\le8$ use the density bounds of DeCorte, Oliveira and Vallentin~\cite{DOV}. These were published with exact verification programs, which we rebuilt and re-ran (Section~\ref{sec:numerics}). The asymptotic bound is the analytic Theorem~1.1 of Bachoc, Passuello and Thiery~\cite{BPT}. Table~\ref{tab:bounds} lists further dimensions.

\begin{theorem}[Formal verification]\label{thm:lean}
The case $n=3$ of Theorem~\ref{thm:A}, in the form of Corollary~\ref{cor:B}, and the deduction ``$m_1(\R^3)\le0.1532996\Rightarrow\chi(\R^3)\ge7$'' are formally verified in Lean~4 with Mathlib. The proof depends only on the standard axioms \texttt{propext}, \texttt{Classical.choice} and \texttt{Quot.sound}.
\end{theorem}

Section~\ref{sec:lean} describes the formal development. It builds on the Lean library accompanying~\cite{OAI}, in particular on its formal proof of the planar rigidity theorem. The only input not formalized is the numerical density bound of~\cite{DOV}, which appears as an explicit hypothesis.

Finally, the same method applied to random independent sets settles a question about finite graphs. For a finite unit-distance graph $G\subset\R^n$ let $\alpha_1(G)=\alpha(G)/|G|$, and set $\alpha_1(\R^n)=\inf_G\alpha_1(G)$. Averaging translates gives $m_1(\R^n)\le\alpha_1(\R^n)$. In the plane, Ambrus et al.~\cite{ACMVZ} and Matolcsi et al.~\cite{MRVZ} conjectured that the inequality is strict, ``an interesting example of a natural infinite graph for which the measurable and the non-measurable variants of a parameter have different values.'' D\'ucz and Varga~\cite{DV} disproved the planar conjecture $\alpha_1(\R^2)=1/4$, and conjectured instead that $f(N)/N=m_1(\R^2)+o(1)$, where $f(N)$ is the least independence number of an $N$-point unit-distance graph in the plane.

\begin{theorem}\label{thm:D}
For every $n\ge2$,
\[
\alpha_1(\R^n)=m_1(\R^n)\quad\text{and}\quad\sup_{G}\chi_f(G)=\sup_G\chi_{gf}(G)=1/m_1(\R^n),
\]
where $G$ ranges over finite unit-distance graphs in $\R^n$ and $\chi_{gf}$ is the geometric fractional chromatic number of~\cite{ACMVZ,MRVZ}. In particular $f(N)=(m_1(\R^2)+o(1))N$, and $\sup_G\chi_f(G)>4.0485$ in the plane.
\end{theorem}

The second equality answers the question of~\cite{MRVZ} whether $\chi_{gf,0}(\R^d)=\chi_{f,0}(\R^d)$ for $d\ge3$. As in~\cite{DV}, the graphs produced are not explicit in practice.

\subsection*{Acknowledgements}
The extension to higher dimensions, Theorem~\ref{thm:D}, and the Lean formalization were developed with the assistance of Claude (Anthropic). The underlying planar results and their formalization are due to~\cite{OAI}.

\section{Setting and inputs}\label{sec:setting}

Let $F=\overline{\Q}\cap\R$ be the field of real algebraic numbers, $E=F(i)\subset\C$, and $K=\{u\in E:|u|=1\}$, a countable divisible abelian group. Write $n=2m$ or $n=2m+1$ and identify $\R^n$ with $\C^m$ or $\C^m\times\R$. The algebraic points form the countable dense subgroup
\[
V=E^m\quad\text{or}\quad V=E^m\times F .
\]
The group $K^m$ acts on $V$ by rotating the $j$-th coordinate plane by the $j$-th component and fixing the last axis. These maps are linear isometries of $\R^n$. The group $G=V\rtimes K^m$ is solvable, hence amenable.

For a unitary representation $U$ of a discrete group $A$ that is a dense subgroup of a Euclidean space (here $E$, $F$ or $V$), let
\[
\cH_c(U)=\{v:\ \|U_zv-v\|\to0\ \text{as }z\to0\text{ in the Euclidean topology}\}.
\]
This is a closed subspace. We use three results of~\cite{OAI}.

\begin{theorem}[{\cite[Theorem 2.3]{OAI}}, Hilbert-space form]\label{thm:rigidity}
Let $U\colon E\to\mathcal U(\cH)$ and $W\colon K\to\mathcal U(\cH)$ be unitary representations with $W_kU_z=U_{kz}W_k$. If $v\perp\cH_c(U)$ and $W_kv=v$ for all $k\in K$, then $\langle v,U_zv\rangle=0$ for every $z\in E\setminus\{0\}$.
\end{theorem}

In the Lean library of~\cite{OAI} this is \texttt{PlaneFiveColor.Spectral.residual\_fourier}. It is deduced there from the rigidity of $K$-invariant probability measures on the dual of $E$ that give no mass to the continuous characters.

\begin{lemma}[{\cite[Lemma 4.1]{OAI}}]\label{lem:factor}
If $U$ is the Koopman representation of a measure-preserving action of $A$ on a probability space, then $\cH_c(U)$ is the $L^2$ space of a factor. Hence its orthogonal projection is a conditional expectation: it preserves $0\le f\le1$ and fixes constants. Translations act continuously on $\cH_c(U)$ and extend to a strongly continuous, positivity-preserving unitary representation of the ambient Euclidean space.
\end{lemma}

The proof in~\cite{OAI} is stated for $A=E$ but uses only that $\cH_c$ is closed, contains constants, and is stable under products of bounded elements and under Lipschitz maps. It applies verbatim to $A=F$ and $A=V$.

\begin{lemma}[{\cite[\S2.2]{OAI}}]\label{lem:mean}
There is a positive, $F$-translation-invariant linear functional $I$ on bounded, boundedly supported functions $F\to\C$ with $I(\varphi|_F)=\int_\R\varphi$ for all $\varphi\in C_c(\R)$. It integrates bounded Hilbert-valued functions weakly, with $\|I(\Phi)\|\le I(\|\Phi\|)$.
\end{lemma}

\section{Two lemmas}

\begin{lemma}\label{lem:comm}
Let $U$ be a unitary representation as above and $Y$ a bounded operator commuting with every $U_z$. Then $Y\cH_c(U)\subseteq\cH_c(U)$. If moreover $Y$ is unitary, or self-adjoint, then $Y\cH_c(U)^\perp\subseteq\cH_c(U)^\perp$, and $Y$ commutes with the orthogonal projection onto $\cH_c(U)$.
\end{lemma}

\begin{proof}
$\|U_zYv-Yv\|=\|Y(U_zv-v)\|\le\|Y\|\,\|U_zv-v\|$. In the unitary or self-adjoint case $Y^*$ also commutes with $U$, so $\langle Yw,h\rangle=\langle w,Y^*h\rangle=0$ for $w\perp\cH_c$ and $h\in\cH_c$. Splitting $v=Pv+(v-Pv)$ gives $PYv=YPv$.
\end{proof}

\begin{lemma}[Twisted rigidity]\label{lem:twisted}
Let $U,W$ be as in Theorem~\ref{thm:rigidity}, and let $Y$ be unitary and commute with every $U_z$ and every $W_k$. If $v\perp\cH_c(U)$ is $W$-fixed, then $\langle v,U_zYv\rangle=0$ for every $z\in E\setminus\{0\}$.
\end{lemma}

\begin{proof}
By Lemma~\ref{lem:comm}, $Yv$ is $W$-fixed and orthogonal to $\cH_c(U)$, and so is $v+sYv$ for every $s\in\C$. Theorem~\ref{thm:rigidity} applied to $v$, $Yv$ and $v+sYv$ gives
\[
0=s\,\langle v,U_zYv\rangle+\bar s\,\langle Yv,U_zv\rangle\qquad(s\in\C).
\]
Taking $s=1$ and $s=i$ shows that both coefficients vanish.
\end{proof}

\section{Proof of Theorem~\ref{thm:A}}\label{sec:proofA}

We prove a slightly more general statement, which is also used in Section~\ref{sec:finite}.

\begin{proposition}[One-label transfer]\label{prop:onelabel}
Let $\mu$ be a $G$-invariant Borel probability measure on a compact metrizable $G$-space $\Omega$. Let $T_a$ $(a\in V)$ and $R_\kappa$ $(\kappa\in K^m)$ be the Koopman operators on $\cH=L^2(\Omega,\mu)$, so that $R_\kappa T_aR_\kappa^{-1}=T_{\kappa a}$. Let $f\in\cH$ satisfy $0\le f\le 1$, $R_\kappa f=f$ for all $\kappa$, and
\[
\langle f,T_af\rangle=0\qquad\text{for every }a\in V\text{ with }|a|=1 .
\]
Let $Q$ be the orthogonal projection onto $\cH_c(T)$, the vectors continuous for the whole $V$-action. Then $p=Qf$ satisfies $0\le p\le 1$, $\E p=\E f$, and $\langle p,T_xp\rangle=0$ for every $x\in\R^n$ with $|x|=1$, where $T_x$ is the extension of Lemma~\ref{lem:factor}.
\end{proposition}

\begin{proof}
For a block $j\le m$ let $U^{(j)}_z=T_{\iota_jz}$, where $\iota_jz$ places $z\in E$ in the $j$-th coordinate plane. Let $W^{(j)}$ be the rotations of that plane, and $P_j$ the projection onto $\cH_c(U^{(j)})$. In odd dimension let $S_t=T_{(0,\dots,0,t)}$ for $t\in F$, and let $P_0$ be the projection onto $\cH_c(S)$.

\emph{Commutation.} All translations commute with one another and with rotations of the other blocks, and axis translations commute with all rotations. By Lemma~\ref{lem:comm}, each $P_j$ (and $P_0$) commutes with every $T_a$, with the rotations of the other blocks, and, through the self-adjoint case of Lemma~\ref{lem:comm}, with the other projections. Covariance gives $\|U^{(j)}_zR_\kappa v-R_\kappa v\|=\|U^{(j)}_{\kappa^{-1}z}v-v\|$, so $P_j$ also commutes with the rotations of its own block. A vector continuous in every block and along the axis is continuous for $V$, so
\[
Q=P_m\cdots P_1\quad\text{or}\quad Q=P_0P_m\cdots P_1 .
\]
Each factor is a conditional expectation (Lemma~\ref{lem:factor}). Hence $0\le p\le1$, and $\E p=\langle f,Q1\rangle=\E f$.

\emph{The plane blocks.} Call a unit vector $a\in V$ \emph{generic} if all of its block components $a_j\in E$ are nonzero. Let $g^{(0)}=f$ and $g^{(j)}=P_jg^{(j-1)}$; these are rotation-fixed. Put $v=(1-P_j)g^{(j-1)}$, which is $W^{(j)}$-fixed and orthogonal to $\cH_c(U^{(j)})$. Write $T_a=U^{(j)}_{a_j}Y$, where $Y=T_{a-\iota_ja_j}$ commutes with $U^{(j)}$ and $W^{(j)}$. Lemma~\ref{lem:twisted} gives $\langle v,T_av\rangle=0$. The cross terms vanish by Lemma~\ref{lem:comm}, so
\[
\langle g^{(j)},T_ag^{(j)}\rangle=\langle g^{(j-1)},T_ag^{(j-1)}\rangle\qquad(a\text{ generic}).
\]
With $g=g^{(m)}$ this gives $\langle g,T_ag\rangle=\langle f,T_af\rangle=0$ for generic $a$.

\emph{Even $n$.} Here $g=p$. The rational points are dense in $S^{n-1}$ (inverse stereographic projection), and the non-generic ones lie in finitely many proper subspheres. So generic algebraic unit vectors are dense, and continuity of $x\mapsto\langle p,T_xp\rangle$ finishes the proof.

\emph{Odd $n$: the axis.} Let $w=g-p\perp\cH_c(S)$. Both $p$ and $w$ lie in the range of $P_m\cdots P_1$, so block translations act continuously on them and extend to real arguments. Since $P_0$ commutes with $T_a$,
\begin{equation}\label{eq:split}
\langle p,T_ap\rangle=-\langle w,T_aw\rangle\qquad(a\text{ generic}).
\end{equation}
Fix an algebraic direction $u$ in the plane blocks with all block components nonzero, and let $\rho(t)=\sqrt{1-t^2}$. For $t\in F\cap(-1,1)$ the vector $a(t)=(\rho(t)u,t)$ is a generic algebraic unit vector, because $F$ is real closed. Let $b(t)=\langle p,T_{(\rho(t)u,t)}p\rangle$; it is continuous in real $t$.

Fix a real $\varphi\in C_c((-1,1))$. The window $\Phi(t)=\varphi(t)\,T_{(\rho(t)u,0)}w$ is continuous with compact support, hence uniformly continuous. Let
\[
\mathcal V=I_t\bigl(S_t\Phi(t)\bigr).
\]
Every $S_t\Phi(t)$ lies in the closed subspace $\cH_c(S)^\perp$, so $\mathcal V$ does too. By translation invariance of $I$,
\[
\|S_s\mathcal V-\mathcal V\|\le I_t\bigl(\|\Phi(t-s)-\Phi(t)\|\bigr)\to0\quad(s\to0),
\]
so $\mathcal V\in\cH_c(S)\cap\cH_c(S)^\perp=\{0\}$. Pairing $\mathcal V$ with $c\,w$ ($c\in\C$) and using~\eqref{eq:split} and Lemma~\ref{lem:mean},
\[
0=I_t\Bigl(\varphi(t)\,\mathrm{Re}\bigl(\bar c\,\langle w,T_{a(t)}w\rangle\bigr)\Bigr)=-\int_\R\varphi(t)\,\mathrm{Re}\bigl(\bar c\,b(t)\bigr)\,dt .
\]
Taking $c\in\{1,i\}$ and applying the fundamental lemma of the calculus of variations gives $b\equiv0$ on $(-1,1)$. Algebraic directions are dense, so by continuity $\langle p,T_xp\rangle=0$ for every real unit vector $x$.
\end{proof}

\begin{proof}[Proof of Theorem~\ref{thm:A}]
Restrict the given coloring to $V$. Let $\Omega\subseteq\{1,\dots,k\}^V$ be the nonempty, closed, $G$-invariant set of proper labelings of the unit-distance graph on $V$, and let $\mu$ be a $G$-invariant probability on $\Omega$, obtained by Følner averaging. Let $f_i=\mathbf 1[\ell(0)=i]$. These are rotation-fixed, satisfy $\sum_if_i=1$, and $\langle f_i,T_af_i\rangle=\mu\{\ell(0)=i=\ell(a)\}=0$ for algebraic unit $a$. Proposition~\ref{prop:onelabel} gives $p_i=Qf_i\in[0,1]$ with $\sum_ip_i=Q1=1$ and zero correlation at every real unit vector.

Choose jointly measurable fields $H_i(x,\omega)\in[0,1]$ representing $T_xp_i$ for almost every $x$, as in~\cite[proof of Proposition 4.2]{OAI}. Then
\[
\int_\Omega\sum_i\int_{B(0,R)}\int_{S^{n-1}}H_i(x,\omega)H_i(x+u,\omega)\,d\sigma\,dx\,d\mu=|B(0,R)|\sum_i\int_{S^{n-1}}\langle p_i,T_up_i\rangle\,d\sigma=0 .
\]
Fix $\omega$ such that all these inner integrals vanish for rational $R$, and such that $\sum_iH_i(\cdot,\omega)=1$ almost everywhere. Color each $x$ by the least $i$ with $H_i(x,\omega)>0$. This coloring is measurable, and its monochromatic unit pairs form a $dx\,d\sigma$-null set.
\end{proof}

\begin{proof}[Proof of Corollary~\ref{cor:B}]
Let $A_i^*$ be the set of Lebesgue density points of $A_i=c^{-1}(i)$. These sets are Borel and pairwise disjoint, and their union is conull by the Lebesgue differentiation theorem. By~\cite[Lemma 4.3]{OAI}, whose proof is dimension-free, $A_i^*$ contains no two points at distance $1$. Since $\sum_i|A_i^*\cap[-T,T]^n|=(2T)^n$ for every $T$, some class has upper density at least $1/k$. Hence $1/k\le m_1(\R^n)$.
\end{proof}

\section{Numerical inputs}\label{sec:numerics}

Table~\ref{tab:bounds} combines Corollary~\ref{cor:B} with published bounds for $m_1(\R^n)$. The rows $3\le n\le 8$ come from~\cite[Table 2]{DOV}. DeCorte, Oliveira and Vallentin provide the dual solutions and an exact verifier (\texttt{rn\_verify}, GMP/MPFR) as ancillary files. We rebuilt the verifier, after one portability fix that makes a base-class destructor virtual and does not affect the arithmetic, and re-ran it on the published data. For $n=3$ it reports ``Solution checked; bound $=0.153299568560593$''. The patched verifier and logs are in the accompanying repository, \url{https://github.com/outsidetext/euclidean-chromatic-transfer}. The rows $n\ge9$ come from~\cite[Table 3]{BPT}. They rest on floating-point linear programs that we have not independently verified; the margin at $n=24$ is about $0.02$. The asymptotic row is the analytic Theorem~1.1 of~\cite{BPT}. Previous lower bounds for $\chi(\R^n)$ are as collected in~\cite{CKR}.

\begin{table}[h]
\centering
\begin{tabular}{rlll}
\toprule
$n$ & $m_1(\R^n)\le$ & previous $\chi(\R^n)\ge$ & new $\chi(\R^n)\ge$\\
\midrule
3 & 0.1532996 & 6 & 7\\
4 & 0.0985701 & 9 & 11\\
5 & 0.0624485 & 9 & 17\\
6 & 0.0450325 & 11 & 23\\
7 & 0.0260782 & 15 & 39\\
8 & 0.0190945 & 19 & 53\\
9--12 & \cite{BPT} & 22, 30, 35, 37 & 69, 90, 116, 164\\
24 & $1.84366\cdot10^{-4}$ & -- & 5424\\
$n\to\infty$ & $(1.268+o(1))^{-n}$ & $(1.239+o(1))^n$ & $(1.268+o(1))^n$\\
\bottomrule
\end{tabular}
\caption{Lower bounds for $\chi(\R^n)$ from Corollary~\ref{cor:B}.}
\label{tab:bounds}
\end{table}

\section{Finite graphs attain the measurable density}\label{sec:finite}

\begin{proof}[Proof of Theorem~\ref{thm:D}]
\emph{(i) $m_1\le\alpha_1$.} Let $A$ be measurable and unit-distance-free, with densities $d_j\to d$ along centered cubes $Q_j$. For a finite unit-distance graph $G$,
\[
\frac1{|Q_j|}\int_{Q_j}|(A-x)\cap G|\,dx=\sum_{g\in G}\frac{|A\cap(Q_j+g)|}{|Q_j|}\to d|G| ,
\]
so some translate of $A$ meets $G$ in an independent set of size at least $(d-\varepsilon)|G|$. (This also shows that the definition of $m_1$ via $\inf_X\sup_A\lambda(A)/\lambda(X)$ used in~\cite{MRVZ} agrees with ours.)

\emph{(ii) $\sup_G\chi_f(G)\ge 1/m_1$.} Suppose every finite subgraph of the unit-distance graph on $V$ has $\chi_f\le w$. Normalizing optimal fractional colorings gives, for each finite $H\subset V$, a random independent set of $H$ covering each vertex with probability at least $1/w$. A weak-$*$ limit in $\{0,1\}^V$, in which cylinder events are clopen, yields a random independent set $I$ of $V$ with $\PP(v\in I)\ge1/w$ for all $v$. Følner averaging over $G$ makes its law invariant without affecting these properties.

Apply Proposition~\ref{prop:onelabel} to $f=\mathbf 1[0\in I]$, and let $H(x,\omega)$ be jointly measurable versions of $T_xp$. For almost every $\omega$, $H(x)H(x+u)=0$ for $dx\,d\sigma$-almost every $(x,u)$. The cube averages $X_T(\omega)=|Q_T|^{-1}\int_{Q_T}H(x,\omega)\,dx\in[0,1]$ have mean $\E p\ge1/w$, so by reverse Fatou $\PP(\limsup_TX_T\ge1/w)>0$. For such an $\omega$, the density points of $\{H(\cdot,\omega)>0\}$ form a measurable unit-distance-free set of upper density at least $1/w$. Hence $w\ge1/m_1$.

\emph{(iii)} The proof of~\cite[Theorem 2]{MRVZ}, $\sup_G\chi_f(G)=1/\alpha_1$, uses only translations through a lattice cube and holds verbatim in $\R^n$. Together with (i) and (ii) this gives $\alpha_1=m_1$.

Averaging a measurable unit-distance-free set over random translates in large cubes and a Haar-random element of $O(n)$ shows $\chi_{gf}(G)\le1/m_1$ for every finite $G$. Since $\chi_f\le\chi_{gf}$, all three suprema coincide.

Finally, $f$ is subadditive (place copies far apart), so Fekete's lemma gives $f(N)/N\to\inf_Nf(N)/N=\alpha_1(\R^2)=m_1(\R^2)$. The value $4.0485$ uses $m_1(\R^2)\le0.24699$~\cite{ACMVZ}.
\end{proof}

\section{Formal verification}\label{sec:lean}

The Lean development, available at \url{https://github.com/outsidetext/euclidean-chromatic-transfer} (about 1{,}800 lines, Lean 4.34.1, Mathlib commit \texttt{d13f23b}) builds on the 70 modules of the library of~\cite{OAI} on which its planar result depends. It models $\R^3$ as $\C\times\R$ with $|(z,t)|^2=|z|^2+t^2$. The main statements are:

\begin{Verbatim}[formatcom=\leanfont]
theorem three_space_transfer {k : ℕ} [NeZero k] (c : ℂ × ℝ → Fin k)
    (hc : ∀ x y : ℂ × ℝ, norm (x.1 - y.1)^2 + (x.2 - y.2)^2 = 1 → c x ≠ c y) :
    ∃ A : Fin k → Set (ℂ × ℝ), (∀ i, MeasurableSet (A i)) ∧
      Pairwise (fun i j => Disjoint (A i) (A j)) ∧
      volume (Set.iUnion A)ᶜ = 0 ∧
      ∀ i, ∀ x ∈ A i, ∀ y ∈ A i, norm (x.1 - y.1)^2 + (x.2 - y.2)^2 ≠ 1

theorem no_six_coloring_of_density (hDOV : DensityBound 0.1532996) :
    ¬ ∃ c : ℂ × ℝ → Fin 6, ∀ x y, usq (x - y) = 1 → c x ≠ c y
\end{Verbatim}

Here \texttt{DensityBound}~$\delta$ states that every measurable unit-distance-free subset of $\C\times\R$ has upper density at most $\delta$ along the cubes $[-T,T]^3$. The modules follow Sections~\ref{sec:setting}--\ref{sec:proofA}:
\begin{itemize}
\item \texttt{Model}: the label space over $E\times F$ and its invariant law.
\item \texttt{Rigidity}: Lemmas~\ref{lem:comm} and~\ref{lem:twisted}, and the plane step.
\item \texttt{Smoothing} and \texttt{LineMain}: the vector-window invariant integral and the axis step.
\item \texttt{Sampling} and \texttt{Density}: measurable fields, the sampled coloring, density points.
\item \texttt{Corollary}: the deduction from \texttt{DensityBound}.
\end{itemize}
\texttt{\#print axioms} reports only \texttt{propext}, \texttt{Classical.choice} and \texttt{Quot.sound}.

\section{Remarks}

\emph{Why the torus suffices.} Theorem~\ref{thm:A} needs neither invariance under all of $SO(n)$ nor a disintegration of spectral measures. A maximal torus $K^m$ of algebraic rotations is abelian, so averaging over it is possible. Rigidity applied to one block at a time, via Lemma~\ref{lem:twisted}, then removes every spectral component that is discontinuous in some block. The remaining axis in odd dimension is handled by the invariant mean, as in~\cite[Lemma 2.5]{OAI}.

\emph{Other forbidden distances.} The argument uses only that the forbidden distance is algebraic. It therefore transfers proper colorings avoiding any finite set of algebraic distances, and gives $\chi\ge1/m_D$ for the corresponding densities.

\emph{The plane.} By~\cite[Theorem 1.3]{OAI}, deciding between $\chi(\R^2)=6$ and $7$ is equivalent to deciding whether weak measurable $6$-colorings of the plane exist. Density and fractional methods cannot settle this, since by Theorem~\ref{thm:D} they are capped at $1/m_1(\R^2)\le 1/0.2293<4.37$.

\begin{thebibliography}{99}
\bibitem{ACMVZ} G.~Ambrus, A.~Csisz\'arik, M.~Matolcsi, D.~Varga, P.~Zs\'amboki, \emph{The density of planar sets avoiding unit distances}, Math.\ Program.\ (2023); arXiv:2207.14179.
\bibitem{BPT} C.~Bachoc, A.~Passuello, A.~Thiery, \emph{The density of sets avoiding distance 1 in Euclidean space}, Discrete Comput.\ Geom.\ 53 (2015); arXiv:1401.6140.
\bibitem{CKR} D.~Cherkashin, A.~Kulikov, A.~Raigorodskii, \emph{On the chromatic numbers of small-dimensional Euclidean spaces}, Discrete Appl.\ Math.\ (2018); arXiv:1512.03472.
\bibitem{Coulson} D.~Coulson, \emph{A 15-colouring of 3-space omitting distance one}, Discrete Math.\ 256 (2002).
\bibitem{DOV} E.~DeCorte, F.~M.~de Oliveira Filho, F.~Vallentin, \emph{Complete positivity and distance-avoiding sets}, Math.\ Program.\ (2022); arXiv:1804.09099.
\bibitem{deGrey} A.~de Grey, \emph{The chromatic number of the plane is at least 5}, Geombinatorics 28 (2018).
\bibitem{DV} \'A.~D\'ucz, D.~Varga, \emph{A unit-distance graph in the plane with independence ratio below 1/4}, arXiv:2606.28157 (2026).
\bibitem{LR} D.~G.~Larman, C.~A.~Rogers, \emph{The realization of distances within sets in Euclidean space}, Mathematika 19 (1972).
\bibitem{MRVZ} M.~Matolcsi, I.~Z.~Ruzsa, D.~Varga, P.~Zs\'amboki, \emph{The fractional chromatic number of the plane is at least 4}, arXiv:2311.10069.
\bibitem{Nechushtan} O.~Nechushtan, \emph{On the space chromatic number}, Discrete Math.\ 256 (2002).
\bibitem{OAI} OpenAI, \emph{The Euclidean plane is not five-colorable}, preprint (2026), \url{https://github.com/openai/math}, with Lean formalization.
\bibitem{OV} F.~M.~de Oliveira Filho, F.~Vallentin, \emph{Fourier analysis, linear programming, and densities of distance avoiding sets in $\R^n$}, J.\ Eur.\ Math.\ Soc.\ 12 (2010).
\bibitem{Raigorodskii} A.~M.~Raigorodskii, \emph{On the chromatic number of a space}, Russian Math.\ Surveys 55 (2000).
\bibitem{RT} R.~Radoi\v{c}i\'c, G.~T\'oth, \emph{Note on the chromatic number of the space}, in Discrete and Computational Geometry, Springer (2003).
\end{thebibliography}

\end{document}
